
\subsection{Overview of the nonlinear-moment landscape}
\label{subsec:nonlinear-landscape-overview}

\paragraph{Why matrix factorization is the testbed.}
The example in \cref{sec:matrix-factorization} should be read as a compact
matrix-factorization instance of a broader tensor-decomposition theme.  Classical
low-rank matrix models represent an object by a small number of latent factors,
while tensor decompositions replace a two-way array by multiway rank-one
components and their variants
\citep{kolda2009tensor,sidiropoulos2017tensor}.  Such tensor decompositions are not
only algebraic objects: in latent-variable learning, observable second- and
third-order moments often reduce parameter recovery to a structured symmetric
tensor decomposition problem
\citep{anandkumar2014tensor}.  This makes matrix and tensor
factorization a natural source of non-convex objectives whose geometry is simple
enough to analyze but rich enough to contain saddles, symmetries, and multiple
factor branches.

The same family of examples has also played a central role in non-convex
optimization theory.  Low-rank matrix recovery and matrix factorization provide
canonical settings where, under appropriate measurement and rank assumptions,
local minima are global or nearly global and saddle points carry negative
curvature
\citep{BhojanapalliNeyshaburSrebro2016,GeJinZheng2017}.  Orthogonal
tensor decomposition gives a parallel strict-saddle model with many symmetry
related optima and saddles
\citep{GeHuangJinYuan2015}.  These works motivate looking not only at first-order
stationarity, but also at the second-variation geometry that decides whether a
stationary point is stable or escapable.

Our population model transfers this matrix/tensor-factorization geometry to a
mean-field product-measure setting.  The two finite factor vectors are replaced
by two mean-field feature maps
\[
    h_{\mu_i}(z)
    =
    \int \sigma(w^\top z)\,\dd\mu_i(w),
    \qquad i=1,2,
\]
and the fitting term compares the matrix-valued products
\[
    h_{\mu_1}(z)h_{\mu_2}(z)^\top
    \quad\text{and}\quad
    h_{\mu^o_1}(z)h_{\mu^o_2}(z)^\top .
\]
Thus the non-convexity is not introduced by a finite-dimensional factor matrix,
but by a nonlinear product map on probability measures.  In the Gaussian-input
spherical specialization, the Hermite expansion in
\cref{lemma:EMN-expansion-F} rewrites this product map in terms of moment tensors
\(T_{i,n}=\int w^{\otimes n}\,\dd\mu_i(w)\).  The fitting loss becomes a squared
residual system for the tensor blocks \(\Delta\mathcal M_p\), while the metric
regularizer controls the Wasserstein distance to the antipodal anchor.  This is
the precise sense in which the appendix studies a nonlinear moment landscape:
the active coordinates are tensor moments, but the actual dynamics remain
Wasserstein dynamics of the underlying measures.

This perspective is close to the mean-field two-layer neural-network viewpoint,
where wide-network training is described by the evolution of a parameter
measure, often through Wasserstein gradient flow
\citep{songmei2018,chizat2018,nitanda2022convex}. It is also close to the
matrix-decomposition example in \citet{yamamoto2025hessianguided}, where a
single mean-field feature map is used through the symmetric self-product
\(h_{\mu}(z)h_{\mu}(z)^\top\).  The present product version separates the two
factor measures and studies the product WGF on
\(\mu_1\otimes\mu_2\).  This separation is useful because the two natural
branches
\[
    \mathcal M_-=\{\mu_2=(-\Id)_\#\mu_1\},
    \qquad
    \mathcal M_+=\{\mu_1=\mu_2\}
\]
become visible in the active odd moments. For
\(n\in\mathcal O\), where \(\mathcal O\) is the finite set of active odd
Hermite degrees, recall
\[
 T_{i,n}=\int w^{\otimes n}\,\dd\mu_i(w),\qquad
 S_n=\frac{T_{1,n}+T_{2,n}}2,\qquad
 A_n=\frac{T_{1,n}-T_{2,n}}2
\]
from \cref{assumption:odd-symmetric-teacher,eq:mf-symmetric-antisymmetric-active-blocks}.
Thus
membership in \(\mathcal M_-\) implies \(\mathbf S=0\), while
membership in \(\mathcal M_+\) implies \(\mathbf A=0\).
These finite moment implications are not converses; see
\eqref{eq:mf-finite-active-moment-map}.
The appendix establishes four local statements: attraction, transverse
descent, and positive second variation in two separate regions. Each statement
has its own hypotheses.

\subsubsection{Summary of the section}

This subsection summarizes the separate population-level statements behind
the diagram in \cref{sec:matrix-factorization}. The local incoming branch,
transverse descent direction, and positive-curvature regions are proved below.
A trajectory connecting these regions, or a finite-sample escape guarantee,
is not established by these geometric statements.
Two conventions apply throughout.
First, throughout this appendix the dynamics are the product Wasserstein gradient
flow, namely the WGF restricted to product measures and product tangent fields, as
recorded in \cref{app:product-measure-wgf}.  This should not be read as a claim
that the unrestricted WGF on all joint measures preserves product structure.
Second, $D^2_{W_2}F(\mu)[v,v]$ is the spherical product pushforward second
variation, defined as the full EXP-pushforward derivative. It equals the
quadratic form of the full
spherical product Hessian $\mathcal H_{\Sphere,\mu}$ in
\cref{lem:mf-full-spherical-Hessian,eq:mf-full-Hessian-second-variation}.
This includes both the measure derivative and the spatial Hessian,
also at stationary states. When the prescribed pushforward path is a product
Wasserstein geodesic, this quantity is its geodesic second derivative; general
admissible fields are not assumed to generate such a geodesic.

The model is fixed by the spherical product objective
\cref{eq:EMN-objective-squared-chi2}.  Under the cap and summability assumptions
in \cref{assumption:on-sphere}, and the finite odd-block symmetric-teacher
assumption in \cref{assumption:odd-symmetric-teacher}, the fitting term admits the
Hermite tensor expansion of \cref{lemma:EMN-expansion-F}. Our Hermite-basis analysis follows a now-standard Gaussian-input toolkit in which activation coefficients directly parameterize optimization\citep{arous2021online,abbe23a, bietti2023} and kernel phenomena~\citep{daniely2016toward,murray2023characterizing}, rather than merely serving as a formal expansion.  The metric regularizer
is
\[
    F_{\mathrm{reg}}(\mu_1,\mu_2)
    =
    \frac{\lambdaReg}{p_W}
    Z_\pm(\mu_1,\mu_2)^{p_W/2},
    \qquad
    p_W=2+\delta_W>2,
\]
where \(Z_\pm\) is the squared product distance to the antipodal Dirac pair
\(\nu_\star^\pm=\frac12\delta_{a_\star}+\frac12\delta_{-a_\star}\).  The
fixed-assignment geometry in \cref{lemma-fixed-assignment} and
\cref{definition-spherical-logarithm-exponential} is used locally on the two caps.
The choice \(p_W>2\) is deliberate: the regularizer gives positive curvature on
metric annuli with \(Z_\pm>0\), but it is Hessian-flat at the antipodal Dirac target
\[
    \mu^\dagger
    \coloneq
    \mu_\star^\pm
    =
    \nu_\star^\pm\otimes\nu_\star^\pm .
\]
Thus the quadratic saddle structure at \(\mu^\dagger\) is governed by the fitting
term.

The finite odd Hermite block gives the active coordinates used to describe the
landscape; their full definition is in
\cref{subsec:introduce-symmetric-antisymmetric-branches-and-saddle-point,eq:mf-symmetric-antisymmetric-active-blocks}.
The measure-level anti-symmetric branch
\(\submanifoldforsaddlepoint\) implies \(\mathbf S=0\) on the active odd block,
whereas the same-density branch \(\mathcal M_+\) implies \(\mathbf A=0\).
These coordinates are diagnostic coordinates for the finite active fitting term,
not a complete parametrization of the full metric-regularized objective.  The
antipodal target belongs to the active-origin intersection and is first-order
stationary, as stated in \cref{prop:intersection-zero-output-fiber}.  The same
subsection also identifies the reference zero-fit teacher point on the
same-density branch in \cref{prop:reference-fit-minimum-on-Mplus}.

\paragraph{Local attraction on the anti-symmetric branch.}
Reflection covariance of the full first variation, proved in
\cref{lem:gradient-symmetry-antisymmetric-branch}, yields branch invariance
up to cap exit in \cref{lem:invariance-antisymmetric-branch}.
The attraction proof in
\cref{subsec:antisymmetric-branch-and-dissipation} then has three steps.
Recall the assigned-pole cost \(c_\star(w)=d_\Sphere(w,\pi_\star(w))^2/2\)
from \cref{eq:mf-fixed-assignment-cost} and its outward gradient
\(q_\star(w)=\nabla_\Sphere c_\star(w)=-\LOG_w\pi_\star(w)\) from
\cref{eq:mf-cost-gradient-LOG}.
The constants \(C_\star\) and \(L_{\mathrm{fit}}\) below depend only on the
fixed activation and teacher; they are defined immediately before the pole
contraction lemma in
\cref{eq:mf-gaussian-derivative-constants,eq:mf-branch-estimate-constants}.
\begin{enumerate}
 \item \textbf{An admissible pole contraction.}
 \Cref{lem:mf-pole-contraction-estimates} uses the spherical transport
 \cref{eq:mf-pole-contraction-path} to show
 \[
 DF_{\mathrm{fit}}(\mu)[(q_\star,q_\star)]\ge-C_\star Z_\pm(\mu)^2,
 \qquad \sup_w\|g_{i,\mathrm{fit}}(\mu,w)\|
                 \le L_{\mathrm{fit}}\sqrt{Z_\pm(\mu)}.
 \]
 Here \(g_{i,\mathrm{fit}}=\nabla_\Sphere(\delta F_{\mathrm{fit}}/\delta\mu_i)\),
 as defined in \cref{eq:mf-branch-radial-notation}.
 \item \textbf{A local radial inequality.}
 Adding the regularizer derivative gives
 \[
 DF(\mu)[(q_\star,q_\star)]
 \ge\lambdaReg Z_\pm(\mu)^{p_W/2}-C_\star Z_\pm(\mu)^2.
 \]
 For \(2<p_W<4\), if
 \(C_\star Z_\pm(\mu)^{(4-p_W)/2}\le\lambdaReg/2\), then
 \cref{lem:radial-positivity-antisymmetric-branch-full-density,eq:mf-full-radial-lower-bound}
 gives
 \[
 DF(\mu)[(q_\star,q_\star)]
 \ge\frac{\lambdaReg}{2}Z_\pm(\mu)^{p_W/2}.
 \]
 The identity
 \[
 \frac{\dd}{\dd t}Z_\pm(\mu_t)=-2DF(\mu_t)[(q_\star,q_\star)]
 \]
 then controls distance while \(\mu_t\in\mathcal M_-\cap
 \mathcal X_\pm(r_{\mathrm{cap}})\).
 \item \textbf{Confinement and convergence from initial data.}
 For \(2<p_W<3\), choose \(z_\star\) by
 \cref{eq:mf-local-basin-threshold}. For every globally well-posed
 characteristic product WGF starting from
 \(\mu_0\in\mathcal M_-\cap\mathcal X_\pm(r_{\mathrm{cap}})\) with
 \(Z_\pm(\mu_0)\le z_\star\), the mean-distance barrier and the pointwise
 fitting-field estimate give cap confinement, dissipation while away from the
 target, and \(\mu_t\to\mu^\dagger\), by
 \cref{prop:reduced-dissipation-antisymmetric-branch}.
 Here globally well posed means existence for all times and uniqueness in the
 characteristic-solution class.
 The distance decay and support barrier are
 \cref{eq:mf-local-distance-decay,eq:mf-local-support-barrier}.
\end{enumerate}
The stationary-state conclusion uses the same direct smallness condition in
\cref{eq:mf-local-radial-smallness}. The distance and support comparisons do
not require a density and apply equally to atomic initial measures.

\paragraph{Spherical product pushforward second variation at anti-symmetric states.}
The positive-curvature estimate is proved in
\cref{subsec:curvature-away-from-saddle-on-antisymmetric}.  The fitting term may
have negative curvature, but on an anti-symmetric active ball
\(\|\mathbf A\|_{\mathcal H_{\mathcal O}}\le R_{\mathrm{act}}\) its negative
second variation is uniformly bounded below by
\cref{lem:fitting-curvature-lower-bound-Mminus-smoothness}.  The metric part gives
the compensating positive curvature: \cref{lem:spherical-squared-distance-Hessian-on-caps}
lower bounds the spherical Hessian of the fixed squared-distance cost on the caps,
and \cref{lem:W2-power-regularizer-Hessian-on-fixed-caps} lifts this to the
regularizer.  Combining the fitting lower bound with the metric lower bound gives
the ambient second-variation bound
\cref{prop:sqchi2-away-from-origin-Wasserstein-convexity}.  Its branchwise
specialization is \cref{cor:branchwise-positivity-Wasserstein-Mminus}.  This is the
appendix-level estimate underlying the annular label in the main-text diagram:
at fixed-assignment states with \(Z_\pm\ge\zeta_{\min}^2\), every admissible
product tangent field has positive spherical product pushforward second variation
with an explicit domination condition
\[
    \lambdaReg\kappa_{\mathrm{cap}}\zeta_{\min}^{\delta_W}
    >
    C_{\mathrm{fit}}^-(R_{\mathrm{act}}) .
\]
Here \(C_{\mathrm{fit}}^-\) is defined in
\eqref{eq:def-Cfit-Mminus-derived}. The field may point outside
\(\mathcal M_-\); global geodesic convexity of the annulus is not asserted.

\paragraph{The escape direction at the saddle.}
The escape direction at the antipodal saddle is constructed in
\cref{subsec:descent-direction-to-escape}.  The teacher imbalance condition is
formulated as a two-sided transverse moment condition in
\cref{assumption:quantitative-unweighted-two-sided-transverse-imbalance}; the
associated transverse moment measure is defined in
\cref{def:two-sided-transverse-teacher-moment-measure}.  The tangential visibility
operator at an antipodal Dirac pair is introduced in
\cref{def:tangential-visibility-matrix-dirac-pair}, and
\cref{lem:active-weighted-transverse-moment-implies-visibility} turns the
transverse teacher signal into a strictly positive escape coefficient.  The
resulting same-density perturbation is summarized in
\cref{lemma:descent-direction-antipodal-dirac-from-unweighted-imbalance}: for a
suitable spherical pushforward path starting at
\(\mu^\dagger=\nu_{\star}^\pm\otimes\nu_{\star}^\pm\) at the fixed anchor
\(a=a_\star\),
\[
    \frac{\dd}{\dd h}\bigg|_{h=0}F(\mu_h)=0,
    \qquad
    F(\mu_h)
    =
    F(\mu^\dagger)
    -
    h^2\gamma_{\mathrm{esc}}
    +
    o(h^2),
    \qquad
    \gamma_{\mathrm{esc}}>0 .
\]
At the same time, because \(Z_\pm(\mu_h)=O(h^2)\) and \(p_W=2+\delta_W\), the
metric regularizer contributes only \(O(h^{2+\delta_W})=o(h^2)\).  Hence the
quadratic escape curvature is not masked by the regularizer. This verifies an available same-density escape direction, not that every
perturbation necessarily follows it.

Our escape argument adapts the zero-output saddle mechanism of
\citet{yamamoto2025hessianguided} to the product-spherical setting by replacing
affine interpolation with a teacher-visible same-density tangent perturbation.

\paragraph{Ambient curvature near the same-density teacher.}
Finally, \cref{subsec:post-escape-convexity-on-the-same-density} proves the
local teacher-side curvature estimate. The branch \(\mathcal M_+\) is
invariant by exchange symmetry, as shown in
\cref{lem:invariance-same-density-branch}. At the same-density teacher point
\(\mu_+^o=\nu^o\otimes\nu^o\), the fitting residual vanishes and its second
variation is a nonnegative square. If
\(Z_o:=Z_\pm(\mu_+^o)>0\), the metric regularizer supplies a strictly positive
margin. The proof of \cref{prop:Mplus-teacher-local-positive-curvature} keeps
this margin positive in a sufficiently small product Wasserstein neighborhood
relative to the fixed-assignment cap class. Its key nearby-state estimate is
\[
 D_{W_2}^2F(\mu)[v,v]
 \ge
 \left[
 \lambdaReg\kappa_{\mathrm{cap}}Z_\pm(\mu)^{(p_W-2)/2}
 -C\|h_{\mu_1}h_{\mu_2}^{\top}
          -h_{\nu^o}h_{\nu^o}^{\top}\|_{L^2_z(\mathrm F)}
 \right]\|v\|_{L^2(\mu;\mathbb R^{d_w})}^2,
\]
where the proof-local \(C<\infty\) depends only on the fixed activation.
The result does not imply that the teacher is stationary for the regularized
objective or that a trajectory escaping the saddle reaches this neighborhood.
